% main.tex
\documentclass{book}

% Puedes añadir paquetes aquí que necesites para todo el documento
\usepackage[utf8]{inputenc}
\usepackage[spanish]{babel}
\usepackage{amsmath} % Ejemplo de paquete
\usepackage{graphicx} % Para imágenes
\usepackage[breakable]{tcolorbox} % Para los cuadros de texto

% Aquí puedes configurar márgenes, etc.
\usepackage[a5paper, margin=1.2cm]{geometry} %para reducir un poco los margenes ya que es un a5

% Fuentes
\usepackage{newcent}    % New Century Schoolbook como fuente principal
\usepackage{courier}     % Courirer disponible para títulos
\usepackage{sectsty}    % Para cambiar la fuente de las secciones

% Aplicar Helvetica solo a los títulos de sección                                                                                                             
\allsectionsfont{\ttfamily}

\setcounter{secnumdepth}{0} %solo numeramos los capítulos

\title{Guía de Glorantha}
\author{wissorl}
\date{\today}

\begin{document}

% Portada
\maketitle

% Tabla de Contenidos, Lista de Figuras, etc. (se generan automáticamente)
\frontmatter % Para numeración romana en elementos preliminares
\tableofcontents
%\listoffigures
%\listoftables
% \include{prologo} % Si tienes un prólogo en archivo separado

\mainmatter % Para numeración arábiga en el contenido principal

% Aquí incluyes tus capítulos individuales
\include{cap0_intro}
% capitulo01_geografia.tex
\chapter{Geografía} % Usas \chapter aquí
\label{cha:geografia} % Opcional: para referencias cruzadas

\textit{Glorantha es un cubo de Tierra rodeado y flotando sobre el interminable Río Sramak. Las aguas, algunas muy profundas, cubren gran parte de la Tierra, dividiendo su superficie en dos continentes y muchas islas. En el centro del cubo se encuentra La Fosa de Magasta, un remolino que desemboca en el Inframundo.}

Sobre la superficie está el Domo Solar, sostenido por las Cuatro Direcciones. En la parte superior del Domo Celestial se encuentra la Estrella Polar; sobre ella está Dayzatar, la Luz Pura. La Luna Roja y Orlanth disputan el Aire Medio entre el Domo Celestial y la Tierra.

Bajo la Tierra y las aguas más profundas se encuentra el Inframundo.

\section{EL MUNDO MUNDANO}
\label{sec:el_mundo_mundano}
El Mundo Mundano, donde existen los seres mortales, se encuentra en la parte superior del cubo de la Tierra. Tiene forma cuadrada, con aproximadamente 8000 kilómetros en cada lado. Está dividido en dos continentes y numerosas islas. Aquí las personas viven, aprenden y mueren, y pueden ascender a otros reinos o volver a entrar en los ciclos de la Naturaleza. En los bordes se encuentran tierras mágicas, y más allá se hallan el vacío, los terrores y el mal del Caos.

El continente del norte es el tema principal de este libro.

\subsection{GENERTELA}
El Continente del Norte del Mundo Mundano se llama Genertela, que significa “Lugar de Genert”, el Rey de la Tierra que una vez gobernó estas tierras. Genertela tiene aproximadamente 5000 km de este a oeste y 1700 km de norte a sur. Al norte está el Glaciar de Valind, una tierra helada gobernada por el dios del invierno, donde no habitan razas mortales. Al sur, este y oeste se extienden los océanos: Homeward, Kahar y Neliomi. Muchos mares y bahías invaden la costa.

Una geografía rudimentaria del continente comienza con los Desiertos de Genert, una vasta región árida que se extiende por todo el continente. El Rey de la Tierra fue asesinado durante la Gran Oscuridad por Wakboth, y nunca regresó para dar vida a su jardín. Esto dividió el Este del Oeste.

Genert engendró muchas hijas sobre la Tierra Abundante, y ellas gobiernan los poderes y riquezas de la Tierra en diferentes regiones del continente. Hay siete hijas que dan nombre a las regiones del continente:
Aquí detallo los antecedentes del estudio.

\begin{enumerate}
  \item  \textbf{Ernalda}
  \item  \textbf{Frona}, que gobierna en el noroeste;
  \item  \textbf{Ralia}, en el centro-oeste;
  \item  \textbf{Seshna}, ahora llamada Madre de las Islas desde su hundimiento;
  \item  \textbf{Pelora}, en el centro-norte;
  \item  \textbf{Teshna}, esposa del dios del sol y amante del planeta rojo; y
  \item  \textbf{Kralora}, Madre de los Hsunchen.
\end{enumerate}

Tres grandes cadenas montañosas, creadas por los dioses en su época, separan a las hijas entre sí. Tres grandes ríos drenan las tierras del centro y el oeste, y muchos otros ríos menores drenan las demás tierras.

\subsubsection{Ernalda}
La diosa de la Tierra, vivía en las estribaciones del Pico, que desapareció antes de que comenzara el Tiempo, pero dejó un considerable Poder en las tierras cercanas. Ernalda es la mayor y más poderosa Diosa de la Tierra en todo el mundo, y su nombre se invoca en tierras mucho más allá de su cuerpo.

Su cuerpo incluye la costa sur de Genertela desde Seshnela hasta Teshnos y hacia el norte más allá de El Paso del Dragón hasta Saird. Las regiones varían ampliamente, desde los densos bosques del oeste de Maniria hasta las exuberantes tierras de cultivo del País Sagrado, desde las colinas y montañas de El Paso del Dragón hasta las tierras desoladas que bordean los Desperdicios de Genert. 
	
El Paso del Dragón es donde el cuerpo de Ernalda cruza la división continental. El monte Kero Fin, cubierto de hielo, una aguja de una montaña de unos diez kilómetros de altura, domina la región. El Paso del Dragón es una de las regiones más militarmente significativas en Genertela, ofreciendo el único paso adecuado para grandes ejércitos a través de las Montañas Rockwood, que se extienden sin interrupción desde aproximadamente 1600 km al oeste y unos 1000 km al este. Su población belicosa y muchas entidades mágicas la convierten en una región intimidante para cualquiera que intente conquistar o pacificar. 

\subsubsection{Frona}
Esta diosa está en el noroeste y está separada de su hermana del sur, Ralia, por montañas. Al norte de su bosque se encuentra el Glaciar de Valind, un páramo helado. El río Janube fluye desde el Mar Dulce al este de ella hasta el Océano Occidental. 

Muchas ciudades-estado salpican el Janube mientras su bosque del norte es hogar de muchos bárbaros primitivos, como los Rathori Bearwalkers y los Reindeer People. El sur montañoso es hogar de muchas tribus bárbaras mantenidas unidas por reyes despiadados. 

Al oeste de Fronela se encuentra la antigua fortaleza Brithini de Sog City, gobernada por hechiceros inmortales. Más allá está el reino Hrestoli de Loskalm, uno de los reinos más poderosos de Glorantha y gobernado por magos meritocráticos. 

\subsubsection{Ralia}
Esta diosa está rodeada de tierra y tiene gran fama entre las Razas Ancianas por darles refugio. Está rodeada por montañas en tres lados, con solo un paso a través de cada lado. La ruta del sur se perdió o se ocultó antes de que comenzara el Tiempo. La mayoría de las tierras cerca de las montañas están rotas y montañosas, dejando el centro y el suroeste bajo y plano. 

Esta parte inferior se centra en el río Tanier y el lago Felster, y a menudo se llama Tanisor o Safelster. Estos nombres se asignaron históricamente a varios reinos o confederaciones allí. 

Las regiones montañosas de Ralios están densamente pobladas con Razas Ancianas, incluidos Aldryami, Trolls, Dragonewts y varias ciudades Mostali, así como muchas tribus de bárbaros primitivos y adoradores de los Portadores de Luz. 	Las regiones bajas han estado bajo control humano desde la Era del Amanecer. Está densamente urbanizada y consta de muchas ciudades-estado enemistadas. 

\begin{tcolorbox} [title=Paso del Dragón,breakable]    

El Paso del Dragón se menciona mucho en este libro. Hablando en sentido estricto, El Paso del Dragón es un paso entre dos cadenas montañosas intransitables que, de otro modo, dividen las vastas llanuras del norte de las costas del sur. Hablando más ampliamente, es el país de colinas circundante que ha sido un cruce cultural desde que comenzó el Tiempo.

El Paso del Dragón es el corazón tradicional de la civilización Orlanthi, una de las primeras civilizaciones humanas que surgieron en la Era del Amanecer. Los Orlanthi reverencian a Orlanth, el Dios de la Tormenta (por quien llevan su nombre), y a Ernalda, la Diosa de la Tierra. Han construido varios imperios en su larga historia, solo para desgarrarlos en conflictos fratricidas. Sus reinos tienden a ser caleidoscopios cambiantes de tribus, gobernados por sacerdotes guerreros.

En el centro de El Paso del Dragón se encuentra el Monte Kero Fin, la madre de Orlanth, el Rey de la Tormenta. Cerca está el Templo Shaker, el gran templo de las diosas de la Tierra protegido por la temida Maran Gor.

    Al sur de Kero Fin se encuentran las colinas y valles de las Grazelands, gobernadas por el Pueblo del Caballo Puro y su Reina del Caballo Emplumado. Los valles están habitados por agricultores que tributan al Pueblo del Caballo Puro y están bajo la protección de la Reina del Caballo Emplumado.
	
    Al oeste y al norte de Kero Fin se encuentra Tarsh, un poderoso reino gobernado desde la ciudad de Furthest por una dinastía lunar leal al Imperio Rojo. Esa parte de Tarsh más cercana a Kero Fin se niega a reconocer a los reyes lunares y se les llama los Exiliados de Tarsh.
	
    Al este de Kero Fin se encuentra el Valle de las Bestias y el Pantano de las Tierras Altas. El Valle de las Bestias está habitado por los Hombres Bestia, híbridos mitad animal, mitad humano, y gobernado por el semidiós Ironhoof el Centauro. El Pantano de las Tierras Altas es una tierra fría y sombría habitada por los muertos vivientes animados por la maligna inteligencia de Delecti el Nigromante.
	
    Más al este se encuentran las veinticuatro tribus del Reino de Sartar, unidas en su capital elevada en Boldhome. La lucha entre Sartar y el Imperio Lunar es la fuente de las Guerras de los Héroes, que están destinadas a poner fin a la Tercera Era de Glorantha.
	
    Al este de Sartar se encuentran las llanuras de Prax, una tierra árida de chaparral y manadas de animales. Está habitada por los nómadas Praxianos. El Río de las Cunas es una delgada franja de tierra cultivable al este de Prax y proporciona alimento para la única ciudad notable en Prax: la Ciudad Libre de Pavis.
	
    Al sur de Sartar se encuentra Heortland, hogar de la confederación tribal Hendriki, centrada en Whitewall. Los colonos y la dinastía gobernante de Sartar provienen de entre los Hendriki.
	
    Al oeste de Heortland se encuentra el Plateau de las Sombras, las ruinas oscuras del palacio del Único Anciano. Es un lugar embrujado, habitado por trolls y fantasmas del pasado.
	
    Al suroeste se encuentra el Reino de Esrolia, el hogar de Ernalda, la Reina de la Tierra. Nochet, la ciudad más rica y poblada de Glorantha, domina esta tierra fértil.
\end{tcolorbox}

\subsubsection{Seshna}
Seshna está rodeada por tres lados de océano. Ella es la madre de los leones Pendali, pero fue colonizada por los Brithini antes del Amanecer. Después del Amanecer, el héroe Hrestol mató a su hija y fue exiliado de Seshnela. Para aplacar a Seshna, el padre de Hrestol, Froalar, se casó con la diosa y sus hijos fueron los Reyes Serpiente. 

Los Reyes Serpiente fundaron un poderoso reino que, en la Segunda Era, se unió a las ciudades-estado de Jrustela para convertirse en el todopoderoso Imperio del Mar Medio. Sus hechiceros dominaron Glorantha y cambiaron el paisaje mítico del cosmos. Pero en su arrogancia, los Aprendices de Dios fueron demasiado lejos y fueron destruidos. Demigods de las Puertas del Oeste navegaron a Seshnela y la destrozaron. Aún vive, pero ahora a menudo se la llama Seshna de las Islas.

\subsubsection{Pelora}
Pelora se encuentra en el centro de varias de sus hermanas y tiene contacto directo con cualquier otro mar excepto a través del peligroso Mar Blanco. Montañas rodean el oeste, sur y este, llenando esas regiones de tierras rocosas y quebradas, formando así el Cuenco Peloriano. 

El Mar Dulce domina las regiones occidentales de Peloria, y de él nace el poderoso río Poralistor, ancho y fértil, que se une al Delta del Trueno en el norte. El río Oslir nace en El Paso del Dragón en las tierras de Ernalda, y también fluye hacia el Delta del Trueno. El poderoso Arcos, tercero de los grandes ríos, fluye desde el sureste. 

Pelora ha sido hogar de muchos imperios, desde el Imperio Solar de la Era de los Dioses hasta el Imperio Dara Happan de la Era del Tiempo. Ahora es hogar del Imperio Lunar.



\begin{tcolorbox} [title=El Imperio Lunar,breakable]   
    
    En el último siglo más o menos, el Paso del Dragón ha caído cada vez más bajo la influencia del Imperio Lunar. Este poderoso imperio surgió de las tierras centrales de la otra gran civilización humana para recuperarse en el Amanecer— la civilización Dara Happan. Dara Happa fue fundada por los Dioses del Fuego durante la Era Dorada, pero apenas sobrevivió a la Guerra de los Dioses. Se recuperó en el Tiempo y ha estado en conflicto con los Orlanthi durante mucho tiempo.
	
    Después de que Dara Happa fuera conquistada y oprimida por invasores extranjeros, siete individuos desesperados recorrieron caminos prohibidos y resucitaron a la diosa de la Luna Roja, asesinada en la Guerra de los Dioses. Esta nueva diosa abrazó el Caos y no estaba sujeta a las limitaciones a las que otros dioses habían acordado someterse para comenzar el Tiempo. Buscaba reclamar su lugar en los cielos, lo que inició una enemistad mítico-mágica con Orlanth, el Dios de la Tormenta. Ella es la diosa del Imperio Lunar y su hijo, el Emperador Rojo, que se reencarna constantemente, gobierna el imperio en su nombre.
	
    La enemistad entre Orlanth y la Luna Roja ahora amenaza con destruir Glorantha. El Paso del Dragón es el principal campo de batalla de esta disputa.
\end{tcolorbox}




\subsubsection{Teshna}
Teshna se encuentra entre Ernalda y Kralora. Es esposa del Sol, pero tomó a su hijo Tolat como amante. Está densamente poblada por humanos, elfos amarillos y enanos. Ha tenido muchos amantes, a veces no del todo por elección, pero sigue siendo rica y exuberante.

\subsubsection{Kralora}
Lejos al este, Kralora ha estado aislada de sus hermanas al oeste en la historia y la continuidad. Es tanto la Madre de los Hsunchen, las tribus animales salvajes, como la devota partidaria del Emperador Dragón. 

El centro de la civilización está alrededor del Suam Chow, un mar poco profundo en la cima de Kralora. La gran Isla del Dragón llamada Hum Chang está al este, y es el hogar del espléndido Imperio del Dragón que gobierna todas las tierras alrededor del Suam Chow.

\subsubsection{Pamaltela}
El Continente del Sur del Mundo Mundano se llama Pamaltela, que significa Lugar de Pamalt, el Rey de la Tierra que aún gobierna esa tierra. Pamaltela mide aproximadamente 6700 km de largo de este a oeste y 3300 km de norte a sur. Al norte está el Océano de Regreso. Al este se encuentra el cálido Mar Togaro, el más antiguo de las grandes aguas. Al sur, más allá del Desierto de Nargan, hay una tierra y mar de fuego interminable, imposible de acercarse. Al oeste, más allá de las tierras de los hombres, se extiende el frío Mar Occidental, que no tiene límites. 

Pamaltela también fue devastada por el Caos, pero se recuperó mejor que Genertela porque Pamalt, el Rey de la Tierra, sobrevivió. La paz relativa y la abundancia continúan en gran parte de la tierra, de modo que muchos humanos viven una vida pastoril exuberante. Muchas Razas Ancianas aún son poderosas: innumerables Aldryami, Mostali, Uz y una variedad de criaturas aisladas y oscuras. Las ciudades humanas salpican las costas del norte.






%\include{capitulo02_metodologia}
%\include{capitulo03_resultados}
% ... y así sucesivamente para cada capítulo

\backmatter % Para apéndices, bibliografía, etc.
% \include{apendiceA} % Si tienes apéndices en archivos separados
\bibliography{referencias} % Si usas BibTeX/BibLaTeX
\end{document}
